\section{Introduction}
Security operations centers (SOCs) face alert overload from heterogeneous monitoring systems. The difficulty is not only volume: operational labels are noisy, alert errors are structurally correlated, and telemetry from hosts, IPs, processes, time windows, and threat-intelligence feeds is connected rather than independent. A useful model must therefore denoise alerts while preserving rare malicious evidence and producing a compact analyst queue.

Existing noisy-label methods, including CoLD-style purification, usually assume that suspicious samples can be filtered as independent records. This assumption is weak in SOC data. Embedding-space similarity alone does not identify whether a label is inconsistent, IID noise assumptions fail when errors propagate through related entities, and hard deletion can remove alerts that are evidentially important for an investigation.

Graph-CoLD addresses these limitations with three mechanisms. First, it introduces Graph-CDM as a label-space consistency diagnostic function over prediction, neighborhood, view, and temporal evidence. Second, it uses evidence-preserving weighting to retain class-frequency and anomaly-supported samples instead of deleting all suspicious alerts. Third, it ranks alerts with a SOC priority score combining predicted maliciousness, structural inconsistency, and evidential strength.
